\chapter{Preliminaries}\label{chapter:preliminaries}

Before proceeding with the approach, it's important to give a precision definition of
\begin{enumerate}
    \item the Zipfian distribution
    \item Recognizing a zipfian distribution
    \item Discrete Variate generation
\end{enumerate}

\section{The Zipfian Distribution}
This distributions has had ample impact in the areas of quantitative linguistics and information theory \cite{powersApplicationsExplanationsZipfs1998a}, as well as within computer science in the fields of storage systems and \ac{RDBMS}\cite{grayQuicklyGeneratingBillionrecord1994}. It states that "(...) when a list when 
a list of measured values is sorted in decreasing order, the value of the $n$-th entry is often approximately inversely proportional to $n$" \cite{ZipfsLaw2025}. Mathematically, this is expressed by

\begin{equation}\label{eqt:zipfs_law}
    \textbf{word frequency} \propto \frac{1}{\textbf{word rank}}
\end{equation}

The form of interest for this research has an additional parameter $\alpha$, yielding

\begin{equation}\label{eqt:zm_law}
    \textbf{frequency} \propto \frac{1}{\textbf{rank}^\alpha}
\end{equation}

The probability distribution for a given rank $k$ given $N$ elements is:

\begin{equation}\label{eqt:zm_prob}
    f(k,N, \alpha) := 
 \begin{cases} 
      \frac{1}{H_{N,\alpha}} \frac{1}{k^\alpha} & 1 \leq k \leq N \\
      0 & k < 1 \, \text{or} \, N < k 
   \end{cases}
\end{equation}

where $H_{N,\alpha}$ is defined by the generalized harmonic number

\begin{equation}\label{eqt:generalized_hn}
    H_{N, \alpha} = \sum_{k = 1}^{N} \frac{1}{k^\alpha}
\end{equation}

To give more concrete form to this distribution, we can visualize how the curve looks like for $N = 100$ and with an
exponent $\alpha = 1$

\begin{figure}
    \centering
    \begin{tikzpicture}
    \begin{axis}[
        xlabel={$k$},
        ylabel={$P(X=k)$},
        grid=both,
        thick,
        ymin=0,
        ymax=0.2,
        xtick distance=10,
    ]
    
    \pgfplotstableread[col sep=comma]{data/zipf_pmf.csv}\datatable  
    \addplot table[x=k, y=PMF] {\datatable};
    \end{axis}
\end{tikzpicture}\caption{Example of a Zipfian distribution with parameters $N,\alpha = (100, 1)$}\label{fig:zipf_pmf}
\end{figure}

\section{Recognizing a Zipfian distribution}
One way to visually recognize whether a sample follows a Zipfian distribution or not is by using a plot in log-log scale.
This is justified by taking equation \ref{eqt:zm_prob} (assuming the probability is always well-defined) and applying
a log on both sides.

\begin{equation}
    \log{f(k,N,\alpha)} = -\alpha\log{k} - \log(H_{N,\alpha})
\end{equation}

Since this is clearly a linear form: $Y = mX + C$, we should expect data that was generated according to a Zipfian distribution
to have a linear form when transformed with the logarithm. Applying this on the last graph yields Figure \ref{fig:loglog_zipf}.

\begin{figure}
    \centering
    \begin{tikzpicture}
    \begin{loglogaxis}[
        xlabel={$k$ (log scale)},
        ylabel={$P(X=k)$ (log scale)},
        grid=both,
        thick
    ]
    
    \pgfplotstableread[col sep=comma]{data/zipf_pmf.csv}\datatable  
    \addplot table[x=k, y=PMF] {\datatable};

    \end{loglogaxis}
\end{tikzpicture}\caption{Application of log on data from \ref{fig:zipf_pmf}}\label{fig:loglog_zipf}
\end{figure}

It's a useful criterium for graphically evaluating if the Zipfian has the expected behaviour through the whole domain. It's not 
foolproof since other power laws also turn linear when transformed with the Zipfian, but since all generators analyse \textit{try} to generate 
such a distribution, the problem disappears.

\section{Discrete Variate Generation}

Theoretically defining a distribution is one aspect of the work, but the main goal is to generate
variates which follow a desired distribution. This can be represented by \ref{eq:rvg_definition}.

\begin{equation}\label{eq:rvg_definition}
F(x) = Pr(X \leq x) \quad -\infty < x < \infty
\end{equation}

The theory for developing variate generators rely on the following assumptions \cite{kuhlHistoryRandomVariate2017}:
\begin{itemize}
    \item There exists a perfect uniform $(0,1)$ generator that produces a sequence of \textit{independent} random variables uniformly distributed in $(0,1)$
    \item Computers can save and manipulate real numbers
\end{itemize}

Which are directly violated once implementation in digital computers is considered. While not \textit{a priori} a problem, since 
computers always introduce different types of errors in numerical computations, it's important to consider by \textit{how much} these
assumptions are violated. Even more so if certain distributions exarcebate problems by this violation \cite{radicchiUnderestimatingExtremeEvents2014}.

There are many possible implementations and theoretical frameworks for creating discrete samplers, be them exact or approximative. Each method has it's
advantages and disadvantages and in the next sections, they will be evaluated. It's important to note, however, that other approaches are possible, but have been omitted since
the ones present in the table have already compared more favourably to them (see \cite{devroyeDiscreteUnivariateDistributions1986,kronmalAliasMethodGenerating1979,kuhlHistoryRandomVariate2017}). Below follows a table (\ref{tab:zipf_sampling}) summarizing the methods which will be analysed later. 
\begin{table}[h]
    \centering
    \renewcommand{\arraystretch}{1.3}
    \setlength{\tabcolsep}{3pt} % Adjust column spacing

    \begin{tabular}{p{3cm} p{3cm} p{4cm} p{5cm}}
        \toprule
        \textbf{Method} & \textbf{Efficiency} & \textbf{Complexity} & \textbf{Best Use Cases} \\
        \midrule
        \textbf{Rejection-Inversion \cite{hormannRejectioninversionGenerateVariates1996}} & Very fast ($O(1)$, best exact method) & More complex, efficient rejection & High-performance simulations, library implementations \\ \hline
        \textbf{Alias Method \cite{kronmalAliasMethodGenerating1979}} & Extremely fast ($O(1)$ after $O(N)$ setup) & Needs precomputed tables, memory-intensive & When support is finite and large $N$ is manageable \\ \hline
        \textbf{Marsaglia’s Layered Rejection \cite{marsagliaFastGenerationDiscrete2004}} & Very fast ($O(1)$, optimized for heavy tails) & Moderate, requires precomputed layer boundaries but low memory use & When many samples are needed efficiently, suitable for Zipf–Mandelbrot with infinite support \\ \hline
        \textbf{Continuous Approximation} & Extremely fast ($O(1)$) & Simple, uses power function & Approximate Zipf-like sampling for large $N$ \\
        \bottomrule
    \end{tabular}
    \caption{Summary of methods for sampling from the Zipf-Mandelbrot distribution}
    \label{tab:zipf_sampling}
\end{table}

